\documentclass[14pt,a4paper]{extarticle}

\usepackage{fontspec}
\usepackage{polyglossia}
\setmainlanguage{russian}
\setotherlanguage{english}

\defaultfontfeatures{Ligatures=TeX}
\setmainfont{Times New Roman}
\newfontfamily\arialfont{Arial}

\usepackage[
    a4paper,
    left=3cm,
    right=1cm,
    top=2cm,
    bottom=2cm,
    headheight=0.7cm,
    headsep=0.4cm,
    footskip=0.7cm
]{geometry}

\usepackage{setspace}
\onehalfspacing

\usepackage{indentfirst}
\setlength{\parindent}{1.25cm}
\setlength{\parskip}{0pt}

\usepackage{ragged2e}
\justifying

\usepackage{titlesec}
\titleformat{\section}
    {\normalfont\arialfont\bfseries\fontsize{14pt}{16pt}\selectfont}
    {\thesection}
    {1em}
    {}
\titleformat{\subsection}
    {\normalfont\arialfont\bfseries\fontsize{14pt}{16pt}\selectfont}
    {\thesubsection}
    {1em}
    {}
\titlespacing*{\section}{0pt}{0pt}{6pt}
\titlespacing*{\subsection}{0pt}{0pt}{6pt}

\usepackage{tocloft}
\renewcommand{\cfttoctitlefont}{\hfill\arialfont\bfseries\fontsize{14pt}{16pt}\selectfont}
\renewcommand{\cftaftertoctitle}{\hfill}
\renewcommand{\cftsecfont}{\normalfont}
\renewcommand{\cftsecpagefont}{\normalfont}
\renewcommand{\cftsubsecfont}{\normalfont}
\renewcommand{\cftsubsecpagefont}{\normalfont}
\setlength{\cftbeforesecskip}{0pt}
\setlength{\cftbeforesubsecskip}{0pt}

\usepackage{caption}
\captionsetup{
    font={small,stretch=1.5},
    labelfont=normalfont,
    justification=centering,
    singlelinecheck=false,
    skip=6pt
}

\usepackage{float}
\usepackage{needspace}
\usepackage{placeins}
\usepackage{array}
\usepackage{longtable}
\usepackage{booktabs}
\usepackage{graphicx}
\usepackage{fancyhdr}
\usepackage{lastpage}
\usepackage{hyperref}
\hypersetup{
    colorlinks=false,
    pdfborder={0 0 0},
    pdftitle={Разработка системы экологического социального мониторинга EcoSignal},
    pdfauthor={Дзантиев Азамат Аланович}
}

\clubpenalty=10000
\widowpenalty=10000
\displaywidowpenalty=10000
\brokenpenalty=10000
\interlinepenalty=100

\pagestyle{fancy}
\fancyhf{}
\fancyheadoffset{0pt}
\fancyhead[C]{\fontsize{12pt}{14pt}\selectfont\itshape Дзантиев Азамат Аланович, отчёт по курсовой работе}
\fancyfoot[C]{\fontsize{11pt}{13pt}\selectfont\thepage}
\renewcommand{\headrulewidth}{0pt}
\renewcommand{\footrulewidth}{0pt}

\begin{document}

\begin{titlepage}
\thispagestyle{empty}
\begin{center}
{\fontsize{14pt}{16pt}\selectfont
МИНИСТЕРСТВО НАУКИ И ВЫСШЕГО ОБРАЗОВАНИЯ\\
РОССИЙСКОЙ ФЕДЕРАЦИИ

\vspace{0.8cm}

Федеральное государственное автономное образовательное учреждение\\
высшего образования\\
«Московский политехнический университет»

\vspace{0.8cm}

Факультет информационных технологий

\vspace{0.3cm}

Кафедра инфокогнитивных технологий

\vfill

{\arialfont\bfseries ПОЯСНИТЕЛЬНАЯ ЗАПИСКА}\\
к курсовой работе

\vspace{0.8cm}

по дисциплине:\\
«Корпоративные информационные системы»

\vspace{0.8cm}

на тему:\\
«Разработка системы экологического социального мониторинга EcoSignal»

\vfill
}
\end{center}

\begin{flushright}
\begin{minipage}{0.48\textwidth}
Выполнил:\\
студент 2 курса\\
специальности\\
«Разработка и интеграция бизнес-приложений»\\
Дзантиев Азамат Аланович\\
Группа: \underline{\hspace{3.8cm}}

\vspace{0.8cm}

Руководитель:\\
Логачев Максим Сергеевич
\end{minipage}
\end{flushright}

\vfill

\begin{center}
Москва\\
2026
\end{center}
\end{titlepage}

\newpage
\tableofcontents
\newpage

\section*{Введение}
\addcontentsline{toc}{section}{Введение}

В современных условиях вопросы экологической безопасности становятся всё более значимыми для городской среды, органов управления и граждан. Загрязнение территорий, несанкционированные свалки, повреждение зелёных насаждений, загрязнение водных объектов и другие экологические проблемы часто выявляются не только специализированными службами, но и обычными жителями. При этом эффективность реагирования во многом зависит от того, насколько быстро информация о проблеме будет зафиксирована, проверена, передана ответственным организациям и доведена до стадии решения.

Актуальность работы связана с необходимостью цифровизации процесса подачи и обработки экологических обращений. Традиционный способ взаимодействия граждан с ответственными ведомствами часто связан с разрозненными каналами связи, отсутствием единой карты происшествий, сложностью отслеживания статуса обращения и недостаточной прозрачностью дальнейшей обработки. В результате часть сообщений теряется, дублируется или требует длительной ручной проверки.

В рамках курсовой работы был разработан MVP системы экологического социального мониторинга EcoSignal. Приложение предназначено для подачи гражданами экологических обращений с указанием координат, описания проблемы и фотографий, а также для дальнейшей обработки этих обращений сотрудниками ответственных организаций. Система объединяет пользовательскую часть, интерактивную карту, серверную часть, базу данных и механизм автоматической модерации обращений.

Целью работы является разработка веб-приложения EcoSignal, обеспечивающего регистрацию, отображение, автоматическую проверку и маршрутизацию экологических обращений граждан.

Для достижения поставленной цели были решены следующие задачи: изучена предметная область экологического социального мониторинга; определены основные роли пользователей системы; спроектирована структура базы данных; реализована регистрация и авторизация пользователей; разработан механизм создания экологических обращений с координатами и изображениями; реализована интерактивная карта обращений; создана панель сотрудника организации для обработки заявок; реализована автоматическая модерация обращений с использованием AI; выполнена проверка основных сценариев работы приложения.

Объектом работы является процесс регистрации и обработки экологических обращений граждан. Предметом работы является программная система EcoSignal, обеспечивающая цифровое сопровождение этого процесса от подачи обращения до его рассмотрения ответственным сотрудником.

При разработке использовались HTML, CSS и JavaScript для клиентской части, PHP для серверной логики, PostgreSQL для хранения данных, JWT для авторизации, Yandex Maps API для отображения обращений на карте, а также OpenRouter API для автоматической AI-модерации обращений. Разработка и запуск MVP выполнялись в локальной серверной среде XAMPP.

Практическая значимость работы заключается в том, что разработанный MVP демонстрирует возможность создания единой цифровой среды для фиксации экологических проблем, их предварительной проверки, отображения на карте и передачи ответственным организациям. Такой подход может быть использован как основа для дальнейшего развития городской или региональной системы общественного экологического мониторинга.

\newpage
\section{Предметная область}

\subsection{Экологические обращения как инструмент общественного контроля}

Экологический социальный мониторинг представляет собой процесс наблюдения за состоянием окружающей среды с участием граждан, общественных организаций и ответственных ведомств. В отличие от исключительно ведомственного контроля, социальный мониторинг позволяет быстрее выявлять локальные проблемы, которые могут оставаться незамеченными при плановых проверках.

Граждане ежедневно сталкиваются с экологическими нарушениями в городской среде: переполненными контейнерными площадками, незаконными свалками, загрязнением водоёмов, вырубкой деревьев, задымлением, неприятными запахами и другими проблемами. Такие ситуации часто имеют конкретную географическую привязку, поэтому для их фиксации важно указывать не только текстовое описание, но и точные координаты, фотографии и категорию проблемы.

В системе EcoSignal экологическое обращение рассматривается как структурированная запись, включающая описание проблемы, категорию, подкатегорию, координаты, изображения, статус обработки и сведения о пользователе. Такой формат позволяет перевести сообщение гражданина из неформального обращения в данные, пригодные для хранения, анализа, маршрутизации и последующей обработки.

\subsection{Проблемы традиционной обработки экологических сообщений}

При традиционной обработке экологических обращений возникает несколько практических проблем. Во-первых, обращения могут поступать через разные каналы: электронную почту, телефон, формы на сайтах, социальные сети или личные обращения. Это затрудняет централизованное хранение и контроль исполнения. Во-вторых, сообщения часто содержат неполные данные: отсутствуют координаты, фотографии, категория проблемы или контактная информация.

Отдельной проблемой является проверка достоверности обращения. Пользователь может ошибочно выбрать категорию, приложить неподходящее изображение или описать ситуацию недостаточно подробно. Если такая проверка полностью выполняется вручную, возрастает нагрузка на сотрудников, а скорость обработки снижается.

Также важна проблема маршрутизации. Экологические вопросы могут относиться к разным организациям: природоохранным ведомствам, лесным службам, региональным департаментам или муниципальным структурам. Без автоматизированной классификации и назначения ответственного подразделения обращение может долго находиться без движения.

\subsection{Участники процесса экологического социального мониторинга}

В рамках разработанного приложения можно выделить несколько основных участников процесса. Первый участник --- гражданин, который фиксирует экологическую проблему, выбирает категорию, указывает место на карте, добавляет описание и фотографии. Для него важны простота интерфейса, возможность быстро отправить обращение и отслеживать его статус.

Второй участник --- сотрудник ответственной организации. Он получает обращения, назначенные его организации или филиалу, просматривает детали, анализирует приложенные изображения, меняет статус обработки и может вести переписку по конкретному обращению. Для этой роли важны фильтрация, удобная панель заявок и доступ к полной информации по каждому случаю.

Третий участник --- автоматизированный модуль модерации. В EcoSignal он используется для предварительной проверки обращения. Модуль анализирует описание, выбранную категорию, координаты и изображения, после чего принимает решение о подтверждении или отклонении обращения, а также помогает определить приоритет и ответственную организацию.

\subsection{Использование геоданных и фотографий при фиксации нарушений}

Геоданные являются ключевым элементом экологического мониторинга, поскольку большинство экологических проблем связано с конкретной территорией. Отображение обращений на карте позволяет видеть распределение проблем, определять зоны повышенной нагрузки и быстрее ориентироваться в ситуации.

Фотографии повышают достоверность обращения и позволяют сотруднику предварительно оценить характер проблемы без немедленного выезда на место. В EcoSignal пользователь может прикрепить изображения к обращению, а система сохраняет их вместе с основной записью. Это особенно важно для таких категорий, как загрязнение воды, незаконная свалка, повреждение зелёных насаждений или признаки пожара.

В пояснительной записке для данного подраздела целесообразно разместить горизонтальный скриншот интерактивной карты EcoSignal с отмеченными обращениями, а также отдельный скриншот формы создания обращения с выбором координат и загрузкой изображений.

\subsection{Роль автоматической модерации в обработке обращений}

Автоматическая модерация используется для снижения нагрузки на сотрудников и повышения качества входящих данных. В разработанном MVP EcoSignal AI-модуль анализирует обращение и возвращает структурированный результат: статус проверки, признак экологической проблемы, соответствие описания и изображения, рекомендуемую категорию, ответственную организацию, филиал, приоритет и краткое объяснение решения.

Такой подход не заменяет окончательное решение сотрудника, но позволяет выполнить первичную фильтрацию обращений. Например, система может отклонить сообщение без фотографии, обращение с нерелевантным изображением или проблему, не относящуюся к экологической тематике. Если обращение подтверждается, оно может быть направлено в ответственную организацию для дальнейшей обработки.

Для MVP такой механизм особенно полезен, поскольку демонстрирует не только хранение заявок, но и более сложный процесс их интеллектуальной обработки. Это делает EcoSignal ближе к корпоративной информационной системе, в которой данные проходят через несколько этапов: создание, проверку, классификацию, назначение ответственного исполнителя, обработку и закрытие.

\end{document}
